\documentclass[11pt]{letter}

\usepackage[utf8]{inputenc}
\usepackage[T1]{fontenc}
\usepackage[margin=1in]{geometry}
\usepackage{hyperref}
\usepackage{siunitx}
\urlstyle{same}

\signature{Jonathan Mace\\Corresponding author}
\address{Jonathan Mace\\Microsoft Research\\Redmond, WA\\jonathanmace@microsoft.com}

\begin{document}

\begin{letter}{%
Professor A.\ V.\ Metrikine\\
Editor-in-Chief\\
\textit{Journal of Sound and Vibration}}

\opening{Dear Professor Metrikine,}
\small

We are pleased to submit the manuscript ``Modal Analysis of a Fluid-Filled
Viscoelastic Oblate Spheroidal Shell: Implications for the Existence of the
`Brown Note','' for consideration as an article in the
\textit{Journal of Sound and Vibration}. The paper addresses a long-standing
popular claim with a rigorously mechanical answer: we model the human abdomen
as a fluid-filled viscoelastic oblate spheroidal shell, identify its low-order
flexural resonances at \SIrange{4}{10}{\hertz}, and show why whole-body
vibration excites these modes far more effectively than airborne infrasound.
An energy-consistent coupling analysis gives an upper-bound mechanical-to-airborne
disparity of approximately $6.6 \times 10^4$, with Monte Carlo uncertainty
quantification confirming that this conclusion is robust across physiological
variation.

The manuscript fits \textit{JSV} squarely: at its core it is a shell-vibration
and modal-analysis problem with fluid loading, long-wavelength acoustic
coupling, and physically interpretable asymptotics. We believe it will also
interest readers beyond structural acoustics, because it connects analytical
vibration theory to biomechanics and occupational health by reconciling
whole-body vibration observations in the \SIrange{4}{8}{\hertz} band with the
near-complete ineffectiveness of airborne sound at the same frequencies. The
analytical framework has already proved broader than its entertaining motivating
question and appears applicable to related problems in biomechanical acoustics.

Potential reviewers who appear well matched to the manuscript are:\\
Marco Amabili, McGill University (\url{marco.amabili@mcgill.ca});\\
Subhash Rakheja, Concordia University (\url{subhash.rakheja@concordia.ca});\\
Philip L.\ Marston, Washington State University (\url{marston@wsu.edu});\\
Farbod Alijani, TU Delft (\url{f.alijani@tudelft.nl}).
We request no reviewer exclusions.

This manuscript has not been published previously, is not under consideration
elsewhere, and all authors have approved its submission. We hope the paper's
engaging premise is matched by the seriousness of its mechanics, and we believe
it will make a memorable and technically solid contribution to the journal.

\closing{Sincerely,}

\end{letter}
\end{document}